\section{Conclusion}
\label{sec:conclusion}

\subsection{Summary}

We have presented Multi-scale MCMC Option B for redistricting, implemented
in the BISECT CLI as \texttt{bisect state --search multiscale}.
The reusable substrate lives in \texttt{bisect-multiscale}; current production
orchestration is in \texttt{bisect-cli}'s \texttt{run\_multiscale} wrapper.
The key results are:

\begin{enumerate}

  \item \textbf{County adjacency from the run's tract graph.}
  The county adjacency graph required for coarse \recom{} moves is derived
  from the existing tract graph by scanning cross-county tract edges.
  No additional shapefiles or census downloads are required for Option B beyond
  the tract GEOID map already paired with the adjacency data.
  The derivation runs in $O(|E|)$ time ($<1$ ms for Texas).

  \item \textbf{Mixing speed improvement for large $k$.}
  On Texas ($k = 38$, $T = 2000$ steps), \ms{} reduces lag-100
  autocorrelation of the edge-cut statistic from $0.51$ to $0.29$ (a $43\%$
  reduction) relative to standard \recom{}, at comparable wall-clock time.
  This yields a $\approx 1.8\times$ increase in effective sample size per
  unit time.

  \item \textbf{Three resolution options via GEOID prefix matching.}
  The resolution system defines Options A (block group $\to$ tract), B
  (tract $\to$ county, implemented), and C (block group $\to$ county), all
  sharing the \texttt{derive\_partition} helper.
  Options A and C require block-group adjacency data; the runner has paths for
  those cases when the data are supplied.

\end{enumerate}

\subsection{When to Use Multi-scale MCMC}

\ms{} is a useful search method when:
\begin{itemize}
  \item $k \geq 30$ (the large-$k$ mixing problem is acute); or
  \item The state has a strong county-district correspondence (many states
        have constitutional or statutory requirements to respect county
        boundaries); or
  \item Ensemble analysis needs faster exploration than single-scale ReCom and
        can tolerate the current approximate, non-reversible coarse move.
\end{itemize}

Standard \recom{} remains appropriate for:
\begin{itemize}
  \item $k \leq 14$ (NC, WI, OR, etc.) where single-scale mixing is adequate;
  \item Rapid exploratory runs where ESS precision is not critical.
\end{itemize}

\subsection{Parameter Guidelines}

Based on the local $\alpha$-sweep results in Section~\ref{sec:mixing}:
\begin{itemize}
  \item $\alpha = 0.3$ is a reasonable starting default for states with $k = 30$--$45$.
  \item $\alpha = 0.35$ may improve mixing for $k \geq 50$ (CA, $k = 52$), but
        should be checked with recorded traces.
  \item $\alpha < 0.2$ provides minimal benefit over standard \recom{}.
  \item $\alpha > 0.45$ increases rebalance failure rates and degrades
        effective acceptance.
\end{itemize}

\subsection{Legal Context}

\paragraph{County-preservation legal requirements.}
Many state constitutions contain explicit county-preservation requirements,
including California Article~XXI \S2(d), Colorado Article~V \S47, and Florida
Article~III \S21.
The multi-scale algorithm's county-level coarse moves treat counties as atomic
units \emph{during the coarse step}, but the subsequent fine-level rebalance can
split counties at the tract level.

The correct legal interpretation is narrow: multi-scale does not inherently
satisfy county-preservation requirements.
Those constraints must be enforced as separate checks, just as with any other
search method.
The coarse moves may tend to produce plans where district boundaries respect
county lines more often than pure tract-level ReCom, but that tendency should be
quantified empirically for the specific state and redistricting cycle.

\subsection{Future Directions}

Three natural extensions of the current implementation:

\paragraph{Multi-scale Forest ReCom.}
Replacing the fine \recom{} chain with Forest ReCom
\citep{dellaLibera2026forestRecom}
would provide stronger target accounting for the fine-level component.
The Metropolis-Hastings acceptance ratio used in \fr{} is well-defined on
the fine plan space; combining it with coarse moves requires care in
defining the joint acceptance criterion.

\paragraph{Options A and C.}
Implementing block-group adjacency loading would enable Options A and C,
providing finer spatial resolution for states with heterogeneous tract
populations (e.g.\ California, where coastal tracts are very small).

\paragraph{Adaptive $\alpha$.}
The coarsening parameter $\alpha$ could be adapted during the run based on
the rebalance failure rate: decrease $\alpha$ when failures exceed a threshold,
increase when failures are rare.
This would automate the $\alpha$ selection currently done manually.

\subsection{Implementation Reference}

The reusable substrate is in \texttt{crates/bisect-multiscale/}: seed
derivation, hierarchy construction, rebalance, and configuration defaults.
Current CLI orchestration lives in the \texttt{bisect-cli} runner. Its entry
point is \texttt{run\_multiscale}; level construction is handled by the same
runner module, and public flags/configuration are wired through the CLI runner
and argument modules.
The standalone \texttt{MultiScaleChain} API is still a consolidation target
rather than the production entrypoint.

Runs are reproducible from the search parameters, fine/coarse levels, base seed
$b$, and the deterministic step-level seed schedule
(Section~\ref{sec:algorithm}); publication packages should also retain the EC
trace and selected-plan hash.
